%! Advanced Factoring — Unit 4, Lesson 4.2 — Guided Notes (Answer Key)
\documentclass[10pt]{article}
\usepackage{algebra2-article}
\usepackage{algebra2-key}   % pulls in -boxes; do NOT also load -boxes
\newcommand{\TallMath}[1]{$\displaystyle #1\rule[-1.4em]{0pt}{3.2em}$}

\newcommand{\lgrid}[4]{%
  \draw[step=1, linegray!40, very thin] (#1,#3) grid (#2,#4);
  \draw[->, semithick, charcoal] (#1,0)--(#2,0) node[below right, font=\scriptsize]{$x$};
  \draw[->, semithick, charcoal] (0,#3)--(0,#4) node[left, font=\scriptsize]{$y$};}

\begin{document}

\pageheader{Unit 4, Lesson 4.2}{Guided Notes --- Answer Key}
\namedateperiod

\vspace{0.08in}

\begin{objectivebox}
\textbf{By the end of this lesson, I will be able to\dots}
\begin{itemize}\setlength{\itemsep}{2pt}
  \item factor out the \emph{GCF} first, in one and two variables, including a negative leading coefficient;
  \item choose a method by the \emph{number of terms}, and apply the \emph{difference of squares}, \emph{perfect-square trinomial}, and \emph{sum and difference of cubes} identities;
  \item factor a four-term polynomial by \emph{grouping} and recognize \emph{quadratic form};
  \item factor \emph{completely}, recognize a \emph{prime} polynomial, and read a function's \emph{zeros} off its factored form.
\end{itemize}
\end{objectivebox}

\vspace{0.05in}

\begin{vocabbox}
\small
Fill in each term as we build it together.
\vspace{2pt}

\termblanklong{Prime (irreducible) polynomial}
\ansline{A polynomial that cannot be factored further over the integers --- e.g.\ $x^2+9$, $2x^2-9$, $x^2+2x+4$.}
\termblanklong{Factored completely}
\ansline{Written as a product in which \emph{every} factor is prime --- nothing is left to pull apart.}
\termblanklong{Sum \& difference of cubes}
\ansline{$a^3+b^3=(a+b)(a^2-ab+b^2)$; \; $a^3-b^3=(a-b)(a^2+ab+b^2)$.}
\termblanklong{Factoring by grouping}
\ansline{Split four terms into two pairs, GCF each pair, then factor out the binomial they share.}
\termblanklong{Quadratic form}
\ansline{A polynomial like $x^4+bx^2+c$ (exponents $2n$ and $n$) that factors as though $x^2$ were the variable.}
\end{vocabbox}

\begin{hookbox}
In Lesson 4.0 you read this graph of $q(x) = x^4 - 5x^2 + 4$.
\begin{center}
\begin{tikzpicture}[scale=0.45]
  \lgrid{-3}{3}{-3}{6}
  \begin{scope}
    \clip (-3,-3) rectangle (3,6);
    \draw[very thick, forest, domain=-2.45:2.45, samples=120, smooth]
      plot (\x,{(\x)*(\x)*(\x)*(\x) - 5*(\x)*(\x) + 4});
  \end{scope}
  \fill[charcoal] (-2,0) circle (3pt);
  \fill[charcoal] (-1,0) circle (3pt);
  \fill[charcoal] (1,0)  circle (3pt);
  \fill[charcoal] (2,0)  circle (3pt);
  \fill[charcoal] (0,4)  circle (3pt) node[right, font=\scriptsize]{$(0,4)$};
\end{tikzpicture}
\end{center}
\begin{itemize}\setlength{\itemsep}{1pt}
  \item Read the four $x$-intercepts off the graph: \ans{$-2,\ -1,\ 1,\ 2$}
  \item Now look at $x^4 - 5x^2 + 4$. Can you \emph{see} those four numbers anywhere in it? \ans{No}
  \item What form of the equation \emph{would} show them? \ansline{The factored form --- a product of linear factors, one per zero.}
\end{itemize}
Yesterday you turned a \textbf{factored} form into a \textbf{standard} form by multiplying. Today you
run the machine backward --- and the prize is the \emph{zeros}. One warning before we start: write
$x^4-16=(x^2-4)(x^2+4)$. Is that factored? \ans{Yes} \; Is it factored \emph{completely}?
\ans{No} \; Why not? \ansline{Because $x^2-4$ is itself a difference of squares: $(x-2)(x+2)$.}
\end{hookbox}

\begin{notesbox}{1. Always Start with the GCF}
The \textbf{GCF} of a polynomial uses the largest common coefficient and the \ans{lowest} power of
every variable the terms \emph{share}.

\vspace{4pt}
\begin{itemize}[leftmargin=1.5em, itemsep=4pt]
  \item $12x^5 - 18x^4 + 30x^3 = $ \ans{$6x^3(2x^2 - 3x + 5)$}
  \item Two variables: \; $15x^3y^2 - 25x^2y^3 = $ \ans{$5x^2y^2(3x - 5y)$}
  \item Negative lead --- factor out $-4x$: \; $-4x^3 + 8x^2 - 12x = $ \ans{$-4x(x^2 - 2x + 3)$}

        Pulling out a \emph{negative} changes the sign of \ans{every term} inside the parentheses.
\end{itemize}

\vspace{4pt}
\textbf{Why first?} Because the GCF can \emph{uncover} a pattern that was hidden. On its own,
$2x^3 - 50x$ matches nothing. Take out $2x$: \; $2x^3 - 50x = $ \ans{$2x(x^2-25)$} \; --- and now the
inside is a \ans{difference of squares}, so the complete answer is \ans{$2x(x-5)(x+5)$}
\end{notesbox}

\begin{notesbox}{2. Count the Terms --- That Tells You Which Tool}
\begin{center}
\renewcommand{\arraystretch}{1.45}
\begin{tabular}{|c|p{10.6cm}|}
\hline
\textbf{Terms} & \textbf{Try this} \\ \hline
$2$ & difference of squares $a^2-b^2$ \; \emph{or} \; sum/difference of \textbf{cubes} $a^3\pm b^3$ \\ \hline
$3$ & perfect-square trinomial, ordinary trinomial factoring, or \textbf{quadratic form} \\ \hline
$4$ & factoring by \textbf{grouping} \\ \hline
\end{tabular}
\end{center}
And the rule that beats all of them: after every step, \ans{look at every factor again.}
\end{notesbox}

\begin{notesbox}{3. Two Terms: Squares}
\textbf{Difference of squares:} \; $a^2 - b^2 = $ \ans{$(a+b)(a-b)$}
\begin{itemize}[leftmargin=1.5em, itemsep=3pt]
  \item $49x^2 - 64 = $ \ans{$(7x+8)(7x-8)$} \qquad ($a = $ \ans{$7x$}, $b = $ \ans{$8$})
  \item Two variables: \; $9x^2 - 25y^2 = $ \ans{$(3x+5y)(3x-5y)$}
  \item Go twice: \; $x^4 - 16 = ($\ans{$x^2-4$}$)($\ans{$x^2+4$}$) = $ \ans{$(x-2)(x+2)(x^2+4)$}
\end{itemize}

\vspace{4pt}
\textbf{The one that will not factor.} A \emph{sum} of squares, $a^2 + b^2$, is
\ans{prime} over the integers. For $x^2 + 4$ we would need two integers whose product is $4$ and
whose sum is \ans{$0$} --- and there are none. So we stop.

\vspace{5pt}
\textbf{Perfect-square trinomials} (three terms, but they belong with the squares):
$a^2 \pm 2ab + b^2 = $ \ans{$(a\pm b)^2$}
\begin{itemize}[leftmargin=1.5em, itemsep=3pt]
  \item $x^2 - 14x + 49 = $ \ans{$(x-7)^2$} \quad Check the middle: $2(x)(7) = $ \ans{$14x$} \checkmark
  \item $4x^2 + 20xy + 25y^2 = $ \ans{$(2x+5y)^2$}
\end{itemize}
\end{notesbox}

\begin{notesbox}{4. Four Terms: Grouping}
Split into two pairs, take the GCF out of \emph{each} pair, then factor out the shared
\textbf{binomial}.

\vspace{4pt}
\textbf{Example.} $x^3 + 5x^2 - 4x - 20$

\vspace{2pt}
\emph{Step 1 --- group in pairs:} \; $(x^3 + 5x^2) + (-4x - 20)$

\vspace{2pt}
\emph{Step 2 --- GCF each pair:} \; \ans{$x^2$}$(x+5)$ \; \ans{$-4$}$(x+5)$

\vspace{2pt}
\emph{Step 3 --- factor out the common binomial:} \ans{$(x+5)(x^2-4)$}

\vspace{2pt}
\emph{Step 4 --- keep going!} \; $x^2 - 4$ is a \ans{difference of squares}, so the complete answer is
\ans{$(x+5)(x-2)(x+2)$}

\vspace{5pt}
\textbf{The sign move.} When the third term is negative, factor a \textbf{negative} out of the second
pair so the two binomials \emph{match}. If your two binomials come out different, the reason is
almost always \ans{a sign}.

\vspace{4pt}
\textbf{Sometimes ``completely'' means stopping.} \; $2x^3 - 6x^2 - 9x + 27 = $ \ans{$(x-3)(2x^2-9)$}

\vspace{2pt}
Can $2x^2 - 9$ be factored over the integers? \ans{No} \; Why not? \ansline{It is not a difference of squares --- $9$ is a perfect square but $2x^2$ is not, and there is no integer pair to use.}
\end{notesbox}

\begin{notesbox}{5. Two Terms: Cubes --- Today's New Identity}
Look back at Warm-Up item 3. You multiplied $(x+2)(x^2-2x+4)$ and got \ans{$x^3+8$}. You proved
today's formula before you knew it had a name.
\[
  a^3 + b^3 = (a+b)(a^2 - ab + b^2) \qquad\qquad a^3 - b^3 = (a-b)(a^2 + ab + b^2)
\]
\textbf{SOAP} --- the sign pattern shared by both identities:
\begin{itemize}[leftmargin=1.5em, itemsep=2pt]
  \item \textbf{S}ame --- the binomial factor takes the \ans{same} sign as the original.
  \item \textbf{O}pposite --- the trinomial's middle term takes the \ans{opposite} sign.
  \item \textbf{A}lways \textbf{P}ositive --- the trinomial's last term is \ans{positive}.
\end{itemize}

\vspace{4pt}
\textbf{Two traps.} The middle term is $ab$, \emph{not} $2ab$ --- this is \ans{never} a
perfect-square trinomial. And the trinomial factor $a^2 \mp ab + b^2$ is always \ans{prime}, so
never try to break it apart.

\vspace{5pt}
Identify $a$ and $b$ \emph{first}, every time.
\begin{itemize}[leftmargin=1.5em, itemsep=4pt]
  \item $x^3 + 8$: \quad $a = $ \ans{$x$}, $b = $ \ans{$2$} \quad $\Rightarrow$ \ans{$(x+2)(x^2-2x+4)$}
  \item $27x^3 - 64$: \quad $a = $ \ans{$3x$}, $b = $ \ans{$4$} \quad $\Rightarrow$ \ans{$(3x-4)(9x^2+12x+16)$}
  \item $8x^3 + 125y^3 = $ \ans{$(2x+5y)(4x^2-10xy+25y^2)$}
  \item GCF first! \; $5x^3 - 40 = 5($\ans{$x^3-8$}$) = $ \ans{$5(x-2)(x^2+2x+4)$}
\end{itemize}
\end{notesbox}

\begin{practicebox}
\textbf{Back to the Hook.} $q(x) = x^4 - 5x^2 + 4$ has three terms --- and the powers $4$ and $2$ are
in \textbf{quadratic form}: it factors as though $x^2$ were the variable.
\begin{enumerate}[leftmargin=1.7em, itemsep=6pt]
  \item Find two numbers whose product is $4$ and whose sum is $-5$: \ans{$-1$} and \ans{$-4$}
  \item So $x^4 - 5x^2 + 4 = ($\ans{$x^2-1$}$)($\ans{$x^2-4$}$)$.
  \item Are you done? Each factor is a \ans{difference of squares}, so keep going:
        $q(x) = $ \ans{$(x-1)(x+1)(x-2)(x+2)$}
  \item A product is zero exactly when one of its factors is zero, so the zeros of $q$ are
        \ans{$1,\ -1,\ 2,\ -2$}
  \item Compare with the $x$-intercepts you read off the graph at the top of these notes. Do they
        match? \ans{Yes}
  \item \textbf{The wall.} Try today's tools on the unit's anchor, $x^3 - 3x + 2$. GCF?
        \ans{No} \; Three terms, but quadratic form? \ans{No} \; Four terms to group?
        \ans{No} \; And yet it \emph{does} factor --- you saw it as $(x-1)^2(x+2)$ in Lesson 4.0.
        What that means: \ansline{Today's tools do not reach every polynomial --- we need a new one (division), which is Lesson 4.3.}
\end{enumerate}
\end{practicebox}

\begin{teachernote}
\textbf{Pacing.} Sections 1--2 are fast (about eight minutes together) --- the GCF is review, and the
term-count table is the organizer students will lean on for the rest of the unit, so leave it on the
board. Spend the time in Sections 4 and 5.

\vspace{3pt}
\textbf{The two answers to insist on.} (1) In Section~3, do not let ``$x^2+4$ is prime'' be an
assertion --- run the integer search out loud until it fails. (2) In Section~5, students who write a
$2$ in the middle term are pattern-matching to $(a\pm b)^2$; the cure is to have them multiply
$(x+2)(x^2-2x+4)$ back out, which they already did in the Warm-Up.

\vspace{3pt}
\textbf{The Practice box is the assessment that matters.} Item~3 separates ``can factor'' from ``can
factor \emph{completely}'' --- students who stop at $(x^2-1)(x^2-4)$ have the method but not the
standard. Item~6 is deliberately a dead end: it is honest about the limits of today's toolkit and
creates the need for Lesson~4.3. If time is short, cut Section~3's two-variable line rather than
item~6.
\end{teachernote}

\end{document}
